\section{Conclusions}

This study pursued four objectives and offers corresponding contributions.
First, the multi-level similarity framework with four hierarchical levels and
ten metrics provides a systematic evaluation of route-matching quality,
revealing that 28.66\% of 1.32 million trip chains achieve exact match with at
least one OTP alternative and that matching degrades sharply with transfer
count (45.6\% for direct trips vs.\ $<$0.3\% for two-or-more transfers). This
framework moves beyond single-metric assessments prevalent in prior work and
quantifies the room for improvement that motivates the subsequent modeling.
Second, the comparison of four model families on an identical hold-out set
demonstrates that theory-grounded models can match or exceed machine learning:
a 24-parameter Latent Class model (67.4\%) outperforms LightGBM (66.2\%), and
a six-parameter MNL achieves 99.7\% of the ML benchmark's accuracy.

Third, the LC analysis reveals that administrative card-type categories are
informative but imperfect behavioral proxies (Cram\'er's V = 0.685): 61\% of
elderly users belong to the Accessibility-Priority class, but 39\% exhibit
Mainstream or Extreme preferences. The elderly subgroup's positive ride-time
valuation and near-lexicographic transfer aversion, confirmed across all four
model families, challenge the assumption that travel time minimization is
universal. Fourth, MNL-based re-ranking of OTP alternatives delivers +1.89\%p
improvement---roughly double the structural ceiling of $\theta$-based
calibration---without re-running the router, with a 10.4:1 win ratio when the
two disagree. The gain reaches +3.17\%p in complex scenarios with five or more
alternatives.

For transit agencies and app developers, MNL re-ranking represents a
lightweight, high-impact upgrade deployable as a post-processing step with no
changes to routing infrastructure. For the Accessibility-Priority segment
(8.7\%), systems should prioritize transfer-free subway options; for disabled
users (walking weight 30.2$\times$, transfer penalty 129.6 min), minimizing
walking distance should take precedence over minimizing travel time. More
broadly, card-type information should serve as a Bayesian prior for
personalization, updated with revealed-preference data rather than treated as
a static label.

Several limitations should be noted. The analysis draws on a single weekday,
though weekday transit demand in Seoul exhibits high day-to-day regularity and
1.32 million chains far exceed typical survey-based studies. Fare
differentials, real-time crowding, and weather are absent but primarily affect
mode choice rather than route choice. The positive $\beta_{\mathrm{ride}}$ for
elderly users invites complementary stated-preference research to disentangle
seat availability from ride-time preference. Future work should extend
re-ranking to full probabilistic route assignment for equity-aware service
analysis, pursue real-time API deployment, multi-city replication,
user-type-specific re-ranking leveraging LC class probabilities, and adaptive
learning algorithms that update taste parameters as trip histories accumulate.
